\section{Применение МКВ в рамках будущего исследования}
\label{sec:future-research}
\setcounter{table}{3}

\subsection{Постановка задачи и сценарий использования}
\label{subsec:mvp-scenario}

Сценарий предусматривает совместное обучение логистической регрессии для прогнозирования прекращения клиентом использования услуги. Три организации предоставят записи собственных клиентов с одинаковым набором числовых признаков и меткой класса. Исходные данные не должны раскрываться другим организациям и координатору.

Предусматриваются два независимых блока: модель МО и подключаемый МКВ-модуль защищённого исполнения. МКВ-модуль будет исполнять модель на трёх серверах, по одному от каждой организации, без включения логики протокола в модель. Координатор будет запускать задания без доступа к исходным данным.

Модель предполагается обучить на синтетических данных, разделённых между организациями, и проверить инференс на отдельной тестовой выборке. Открытое исполнение той же модели послужит базой сопоставления предсказаний и численных отклонений. Результат инференса будет доступен только владельцу запроса.

Цель MVP --- продемонстрировать обучение и инференс через заменяемый МКВ-модуль без раскрытия входов в рамках semi-honest модели при честном большинстве (см. раздел~\ref{sec:theoretical-foundations}). Демонстрация будет подтверждать работоспособность интеграции, но не заменять доказательство безопасности выбранного протокола.

Проверка максимальной точности, устойчивости к активному противнику, fairness и работы под промышленной нагрузкой в MVP не предусматривается. Выбор логистической регрессии и готовых библиотечных операций будет обусловлен простотой реализации и воспроизводимостью, а не максимальной эффективностью или защищённостью. Прототип не предполагается использовать как промышленное решение.

\subsection{Архитектура продукта}
\label{subsec:product-architecture}

Продукт предполагается разделить на блок МО, задающий модель и порядок обучения, и МКВ-модуль защищённого исполнения. Подготовка данных и координация заданий будут обслуживающими функциями, а не самостоятельными вычислительными блоками. Разделение позволит проверять модель в открытом режиме и подключать защищённое исполнение через согласованный интерфейс. Приоритетом MVP будет простота интеграции, а не универсальная поддержка моделей.

Блок МО будет содержать описание логистической регрессии, порядок признаков и правила обновления весов. Вход обучения --- матрица числовых признаков и вектор бинарных меток; вход инференса --- запись того же формата без метки. Выходом станет оценка вероятности прекращения использования услуги. Владелец запроса преобразует её в класс по установленному порогу после получения результата. Для упрощения MVP это исключит отдельное защищённое сравнение при выдаче ответа.

МКВ-модуль будет выполнять операции над защищёнными значениями на трёх серверах. Описание модели останется публичным; входы, метки, веса и промежуточные результаты будут представлены долями. Предварительное вычисление открытого предсказания с последующим разделением результата не допускается. Криптографические операции будут отделены от описания модели; их основания и сопоставление приведены в разделе~\ref{sec:cryptographic-foundations} и таблице~\ref{tab:mpc-primitives}.

Интерфейс блоков будет включать инициализацию параметров, обучение, инференс и сохранение состояния. Публичная конфигурация задаст версию модели, размерности, параметры квантования, число итераций и получателя результата. Приватные входы и состояние будут передаваться через ссылки на защищённые объекты, без открытых массивов. Модуль вернёт ссылку на обновлённое состояние, статус и результат для разрешённого получателя. Блок МО будет задавать операции, но не управлять обменом и восстановлением значений.

Скрытие данных от модели означает отсутствие открытых значений у кода, описывающего вычисления, а не защиту от модели как самостоятельного противника. Этот код будет обращаться к защищённым объектам через МКВ-модуль. Другие организации не получат исходные записи; каждому серверу будут переданы предусмотренные протоколом доли. Границы доступа приведены в таблице~\ref{tab:mvp-architecture}.

{\small
\setlength{\tabcolsep}{3pt}
\renewcommand{\arraystretch}{1.2}
\setlength{\tablecontentwidth}{\dimexpr\textwidth-6\tabcolsep-4\arrayrulewidth\relax}
\begin{longtable}{|P{0.22\tablecontentwidth}|P{0.39\tablecontentwidth}|P{0.39\tablecontentwidth}|}
\caption{Компоненты продукта и границы раскрытия данных}\label{tab:mvp-architecture}\\
\hline
\multicolumn{1}{|H{0.22\tablecontentwidth}|}{\textbf{Компонент}} & \multicolumn{1}{H{0.39\tablecontentwidth}|}{\textbf{Функция}} & \multicolumn{1}{H{0.39\tablecontentwidth}|}{\textbf{Доступные данные}} \\
\hline
\endfirsthead
\multicolumn{3}{@{}l@{}}{\normalsize Продолжение таблицы~\thetable}\\
\hline
\multicolumn{1}{|H{0.22\tablecontentwidth}|}{\textbf{Компонент}} & \multicolumn{1}{H{0.39\tablecontentwidth}|}{\textbf{Функция}} & \multicolumn{1}{H{0.39\tablecontentwidth}|}{\textbf{Доступные данные}} \\
\hline
\endhead
Клиент организации & Проверка формата, подготовка и квантование входа; передача долей серверам. & Собственные записи и метки; публичная конфигурация задания. \\
\hline
Блок МО & Описание модели, последовательности операций и правил обучения. & Публичное описание; ссылки на защищённые значения, без их раскрытия. \\
\hline
МКВ-модуль: три сервера & Защищённое обучение и инференс; хранение состояния; выдача результата. & Локальные доли входов, весов и промежуточных значений; сообщения протокола. \\
\hline
Координатор & Запуск задания и сбор статусов. & Публичные настройки и служебные сведения; без записей, долей и весов. \\
\hline
Владелец запроса & Восстановление результата инференса и применение порога классификации. & Собственный запрос и его предсказание; без чужих входов и весов. \\
\hline
\end{longtable}
}

Поток данных: клиент $\rightarrow$ локальная подготовка и квантование $\rightarrow$ МКВ-модуль на трёх серверах $\rightarrow$ результат владельцу запроса. Блок МО передаст описание вычислений, координатор --- управляющее задание. Порядок признаков, масштабирование и границы значений будут согласованы до запуска. Заранее заданные параметры нормализации позволят упростить MVP, исключив отдельный протокол расчёта статистик объединённой выборки.

Веса после обучения будут сохраняться распределённо: каждый сервер будет хранить собственную долю и идентификатор конфигурации. Последующий инференс будет использовать это состояние без восстановления полного набора весов у координатора. Открытая контрольная модель будет храниться отдельно и использоваться только в эксперименте с синтетическими данными. Журналы будут содержать статусы, размерности и время выполнения, но не записи, доли, градиенты или предсказания; скрытие служебных сведений в MVP не предусматривается.

Заменяемость модуля потребует совместимости интерфейса, числового представления и правил раскрытия. Для упрощения MVP предусматривается адаптер с фиксированными операциями логистической регрессии, а не преобразователь произвольных сетей. Замена реализации потребует проверки численной совместимости и модели угроз, но не изменения описания модели. Оптимизация обмена и расширение вычислений будут рассматриваться после демонстрации работоспособности с учётом раздела~\ref{sec:limitations}.

\subsection{Модель угроз и границы MVP}
\label{subsec:mvp-threat-model}

Модель угроз прототипа будет конкретизировать постановку из пп.~1.3--1.4 применительно к архитектуре подраздела~\ref{subsec:product-architecture}. Предусматриваются три организации и три вычислительных сервера, по одному от каждой организации. Координатор будет управлять заданиями без доступа к приватным входам и долям. Допущения приведены в таблице~\ref{tab:mvp-threat-model}.

{\small
\setlength{\tabcolsep}{3pt}
\renewcommand{\arraystretch}{1.2}
\setlength{\tablecontentwidth}{\dimexpr\textwidth-4\tabcolsep-3\arrayrulewidth\relax}
\begin{longtable}{|P{0.27\tablecontentwidth}|P{0.73\tablecontentwidth}|}
\caption{Допущения модели угроз MVP}\label{tab:mvp-threat-model}\\
\hline
\multicolumn{1}{|H{0.27\tablecontentwidth}|}{\textbf{Параметр}} & \multicolumn{1}{H{0.73\tablecontentwidth}|}{\textbf{Принятое допущение}} \\
\hline
\endfirsthead
\multicolumn{2}{@{}l@{}}{\normalsize Продолжение таблицы~\thetable}\\
\hline
\multicolumn{1}{|H{0.27\tablecontentwidth}|}{\textbf{Параметр}} & \multicolumn{1}{H{0.73\tablecontentwidth}|}{\textbf{Принятое допущение}} \\
\hline
\endhead
Противник & semi-honest: наблюдает состояние своего сервера и сообщения, но не изменяет предписанные операции. \\
\hline
Допустимый сговор & Не более одной организации (её клиент и сервер) за весь цикл работы. Допускается сговор с координатором. \\
\hline
Честное большинство & Не менее двух серверов честны. Административные области независимы; общий доступ к состояниям серверов исключён. \\
\hline
Каналы и среда & Аутентифицированные конфиденциальные каналы; корректная реализация и случайность. Компрометация среды и побочные каналы исключены из модели. \\
\hline
Допустимое раскрытие & Конфигурация, размеры выборок, время и статусы. Предсказание раскрывается владельцу запроса; веса и промежуточные значения не восстанавливаются. \\
\hline
\end{longtable}
}

Защите будут подлежать чужие записи, метки и промежуточные значения; собственные входы противнику уже известны. Допустимость сговора является требованием к протоколу, а не следствием наличия трёх серверов. Зависимости от вспомогательных компонентов подготовки вычислений будут учитываться при выборе платформы. Три процесса на одном компьютере позволят проверить интеграцию, но не независимость административных областей.

Защита от активного противника, проверка истинности данных, предотвращение подмены долей и намеренного прекращения участия не входят в MVP. Выбор semi-honest модели позволит упростить реализацию, исключив проверки активной безопасности, а не обеспечить максимальную защищённость. При отказе сервера задание завершится с ошибкой; устойчивость к отказам, гарантированная выдача результата и fairness не заявляются.

MVP будет ограничен логистической регрессией и синтетическими данными, без поддержки сетей произвольной глубины, промышленной нагрузки и автоматического масштабирования. Вывод сведений об обучающей выборке из разрешённых предсказаний, в том числе при повторных запросах, остаётся вне заявленной защиты. DP и скрытие служебных сведений не предусматриваются; эти ограничения рассмотрены в разделе~\ref{sec:limitations}. Гарантии за пределами принятых допущений не заявляются.

\subsection{Выбор протокола и платформы}
\label{subsec:mvp-protocol-platform}

Согласно таблице~\ref{tab:mpc-primitives}, для арифметической части модели предполагается использовать разделение секрета. В качестве базового сочетания выбираются TFEncrypted и semi-honest вариант ABY3 для трёх серверов \cite{20,21}. Выбор будет обусловлен доступностью библиотечных операций и соответствием допущениям подраздела~\ref{subsec:mvp-threat-model}, а не максимальной производительностью или защищённостью. Обзор применений платформ приведён в подразделе~\ref{subsec:ml-applications}; их сопоставление для данного MVP представлено в таблице~\ref{tab:mvp-platforms}.

{\small
\setlength{\tabcolsep}{3pt}
\renewcommand{\arraystretch}{1.2}
\setlength{\tablecontentwidth}{\dimexpr\textwidth-6\tabcolsep-4\arrayrulewidth\relax}
\begin{longtable}{|P{0.19\tablecontentwidth}|P{0.39\tablecontentwidth}|P{0.42\tablecontentwidth}|}
\caption{Сопоставление платформ для реализации MVP}\label{tab:mvp-platforms}\\
\hline
\multicolumn{1}{|H{0.19\tablecontentwidth}|}{\textbf{Платформа}} & \multicolumn{1}{H{0.39\tablecontentwidth}|}{\textbf{Основание применения}} & \multicolumn{1}{H{0.42\tablecontentwidth}|}{\textbf{Решение для MVP}} \\
\hline
\endfirsthead
\multicolumn{3}{@{}l@{}}{\normalsize Продолжение таблицы~\thetable}\\
\hline
\multicolumn{1}{|H{0.19\tablecontentwidth}|}{\textbf{Платформа}} & \multicolumn{1}{H{0.39\tablecontentwidth}|}{\textbf{Основание применения}} & \multicolumn{1}{H{0.42\tablecontentwidth}|}{\textbf{Решение для MVP}} \\
\hline
\endhead
CrypTen \cite{22} & Интерфейс на основе PyTorch; арифметические и бинарные доли. & Не выбран: режим с доверенным генератором троек добавит допущение. Репозиторий архивирован. \\
\hline
TFEncrypted \cite{21} & Интерфейс TensorFlow; ABY3 на трёх серверах; операции обучения и инференса. & Выбран: готовые операции соответствуют составу участников. Потребуется фиксация совместимых версий. \\
\hline
SecureML \cite{23} & Исследовательская реализация МО на двух не вступающих в сговор серверах. & Не выбран: перенос в заданную трёхсерверную архитектуру потребует дополнительных решений. \\
\hline
\end{longtable}
}

Для линейных операций будет использоваться реплицированное аддитивное разделение в ABY3, без самостоятельной реализации схемы Шамира \cite{20,21}. Сложение, умножение матриц и усечение фиксированной точки будут выполняться библиотекой. Применение готового варианта сократит число собственных криптографических процедур; переход к иной схеме не требуется для демонстрации обучения и инференса. Представление чисел и допустимые отклонения будут зафиксированы в требованиях к прототипу.

Для сигмоиды выбирается библиотечная аппроксимация с тремя линейными участками, предусмотренная реализацией ABY3 \cite{21}. Это проще разработки собственного алгоритма нелинейных вычислений. Выбор участка будет выполняться защищёнными сравнениями над бинарными долями и библиотечными преобразованиями представлений, без раскрытия результатов сравнений. Интеграция отдельной подсистемы GC не предусматривается; используемый библиотекой OT не будет реализовываться заново. Такая конфигурация сохраняет нелинейную обработку внутри МКВ-модуля, а не переносит её на открытые значения.

Для обучения предполагается фиксированное число шагов обновления весов без защищённого критерия остановки. Открытая контрольная реализация будет применять ту же аппроксимацию; её влияние на качество будет дополнительно оцениваться относительно исходной сигмоиды. Это позволит отделить отклонения квантования и защищённого исполнения от изменения самой модели. Пороговая классификация останется локальной операцией владельца результата согласно подразделу~\ref{subsec:product-architecture}.

FHE, самостоятельная интеграция HE и GC, а также протоколы защиты от активного противника не выбираются для MVP. Они потребуют дополнительных механизмов и настроек, избыточных относительно проверки работоспособности. Это решение следует ограничениям таблицы~\ref{tab:mpc-primitives} и принятой модели угроз, а не утверждает непригодность этих подходов для других задач.

Подготовка случайности будет выполняться средствами выбранного трёхстороннего протокола, без отдельного доверенного генератора троек и без передачи этой роли координатору \cite{20,21}. Для воспроизводимости будет зафиксирована версия TFEncrypted с проверенными операциями; наличие криптографической случайности, раздельное размещение серверов и защищённые каналы будут проверяться отдельно. TFEncrypted позиционируется как экспериментальное программное обеспечение \cite{21}; его выбор не заменяет проверки реализации. Адаптер из подраздела~\ref{subsec:product-architecture} позволит заменить платформу без изменения описания модели.

\subsection{Функциональные и нефункциональные требования}
\label{subsec:mvp-requirements}

Требования будут определять проверяемый объём прототипа в рамках подразделов~\ref{subsec:mvp-scenario}--\ref{subsec:mvp-protocol-platform}. Их выполнение должно подтвердить обучение и инференс через отдельный МКВ-модуль, а не достижение промышленной производительности. Функциональные требования приведены в таблице~\ref{tab:mvp-functional-requirements}; форматы и параметры будут фиксироваться общей конфигурацией задания.

{\small
\setlength{\tabcolsep}{3pt}
\renewcommand{\arraystretch}{1.2}
\setlength{\tablecontentwidth}{\dimexpr\textwidth-4\tabcolsep-3\arrayrulewidth\relax}
\begin{longtable}{|P{0.12\tablecontentwidth}|P{0.88\tablecontentwidth}|}
\caption{Функциональные требования к MVP}\label{tab:mvp-functional-requirements}\\
\hline
\multicolumn{1}{|H{0.12\tablecontentwidth}|}{\textbf{Код}} & \multicolumn{1}{H{0.88\tablecontentwidth}|}{\textbf{Требование}} \\
\hline
\endfirsthead
\multicolumn{2}{@{}l@{}}{\normalsize Продолжение таблицы~\thetable}\\
\hline
\multicolumn{1}{|H{0.12\tablecontentwidth}|}{\textbf{Код}} & \multicolumn{1}{H{0.88\tablecontentwidth}|}{\textbf{Требование}} \\
\hline
\endhead
Ф1 & Система должна принимать локальные файлы CSV с согласованными признаками и, при обучении, метками; проверять размерности, диапазоны и отсутствие пропусков до передачи данных. Несоответствующий вход должен отклоняться. \\
\hline
Ф2 & Система должна выполнять квантование и разделение входов на стороне владельца, передавать доли трём серверам и не передавать исходные записи или доли координатору. \\
\hline
Ф3 & Система должна обучать логистическую регрессию через МКВ-модуль с заданным числом обновлений весов и библиотечной аппроксимацией сигмоиды, без включения криптографической логики в описание модели. \\
\hline
Ф4 & Система должна сохранять и загружать состояние модели как локальные доли на серверах с идентификатором конфигурации, без восстановления полного набора весов. \\
\hline
Ф5 & Система должна выполнять защищённый инференс и восстанавливать оценку вероятности только у владельца запроса. Преобразование в класс должно выполняться локально после выдачи результата. \\
\hline
Ф6 & Система должна обеспечивать открытое контрольное исполнение на синтетических данных и сопоставление с разрешёнными результатами защищённого режима; формировать отчёт о качестве, времени и трафике. \\
\hline
\end{longtable}
}

Нефункциональные требования представлены в таблице~\ref{tab:mvp-nonfunctional-requirements}. Значения являются проектными целями, а не результатами измерений. Для упрощения MVP будет использоваться одна конфигурация модели и числового представления, без автоматического подбора параметров.

{\small
\setlength{\tabcolsep}{3pt}
\renewcommand{\arraystretch}{1.2}
\setlength{\tablecontentwidth}{\dimexpr\textwidth-4\tabcolsep-3\arrayrulewidth\relax}
\begin{longtable}{|P{0.25\tablecontentwidth}|P{0.75\tablecontentwidth}|}
\caption{Нефункциональные требования и целевые параметры MVP}\label{tab:mvp-nonfunctional-requirements}\\
\hline
\multicolumn{1}{|H{0.25\tablecontentwidth}|}{\textbf{Показатель}} & \multicolumn{1}{H{0.75\tablecontentwidth}|}{\textbf{Цель и условия проверки}} \\
\hline
\endfirsthead
\multicolumn{2}{@{}l@{}}{\normalsize Продолжение таблицы~\thetable}\\
\hline
\multicolumn{1}{|H{0.25\tablecontentwidth}|}{\textbf{Показатель}} & \multicolumn{1}{H{0.75\tablecontentwidth}|}{\textbf{Цель и условия проверки}} \\
\hline
\endhead
Состав системы & Три организации, три вычислительных сервера и один координатор; модель угроз --- согласно подразделу~\ref{subsec:mvp-threat-model}. \\
\hline
Контрольные данные & 900 обучающих записей, по 300 у каждой организации; 300 отдельных тестовых записей. 20 числовых признаков, бинарная метка; признаки в диапазоне $[-1;1]$. \\
\hline
Числовой формат & Фиксированная точка: 64-битное представление, 16 дробных битов, масштаб $2^{16}$, в соответствии с доступной конфигурацией библиотеки \cite{21}. Усечение --- средствами МКВ-модуля. \\
\hline
Режим обучения & Нулевые начальные веса; шаг обучения $1/16$; 150 обновлений с размером пакета 30 записей, по 10 от каждой организации. Порог классификации --- $0{,}5$. \\
\hline
Численная точность и качество & Снижение доли правильных классификаций --- не более 2 процентных пунктов относительно открытой модели с той же аппроксимацией. Среднее абсолютное отклонение вероятностей --- не более $0{,}005$. \\
\hline
Сетевой обмен & Ориентировочный бюджет --- не более 10~МиБ в среднем на один шаг обучения для указанного пакета. Учитываются исходящие сообщения трёх серверов, включая подготовку, выполняемую внутри шага. \\
\hline
Испытательная среда & Три узла, каждый не менее двух процессорных ядер и 8~ГиБ оперативной памяти; сеть не менее 100~Мбит/с. Время обучения и инференса измеряется без обязательного предельного значения. \\
\hline
Воспроизводимость & Пять повторных защищённых запусков с одинаковыми данными, порядком пакетов и настройками. Фиксируются версии зависимостей и параметры генерации синтетических данных. \\
\hline
\end{longtable}
}

Точность будет измеряться на одной тестовой выборке при одинаковом пороге классификации. Предел в 2 процентных пункта означает абсолютную разность долей правильных ответов, а не относительное снижение на 2~\%. Оба критерия качества должны выполняться в каждом из пяти запусков. Сопоставление с открытой моделью, использующей исходную сигмоиду, будет приводиться отдельно, чтобы учитывать эффект аппроксимации согласно подразделу~\ref{subsec:mvp-protocol-platform}.

Объём обмена будет суммироваться по отправленным сообщениям без повторного учёта на стороне получателей. Первоначальная загрузка входов и настройка соединений будут измеряться отдельно. Превышение бюджета трафика должно отражаться в отчёте, но не требовать перехода к более сложному протоколу ради оптимизации MVP. Параметры среды будут фиксироваться для сопоставимости измерений.

Достаточность разрядности будет проверяться до защищённого запуска для заданных диапазонов и числа шагов, без раскрытия промежуточных значений. Журналы не должны содержать входы, доли, веса, градиенты или предсказания. Генератор синтетических данных и криптографическая случайность будут инициализироваться независимо; повторное использование секретной случайности между запусками недопустимо.

\subsection{План реализации и критерии готовности MVP}
\label{subsec:mvp-implementation-plan}

Этапы реализации приведены в таблице~\ref{tab:mvp-implementation-stages}. Для упрощения MVP будут использоваться готовые операции с интерфейсом подраздела~\ref{subsec:product-architecture}, без собственных криптографических процедур. Данные и настройки задаются таблицей~\ref{tab:mvp-nonfunctional-requirements}.

Бинарные метки синтетических данных будут задаваться порогом линейной комбинации признаков; выборки будут генерироваться раздельно. Такая зависимость упростит проверку обучения, но не будет воспроизводить статистику реальных клиентов.

{\small
\setlength{\tabcolsep}{3pt}
\renewcommand{\arraystretch}{1.2}
\setlength{\tablecontentwidth}{\dimexpr\textwidth-6\tabcolsep-4\arrayrulewidth\relax}
\begin{longtable}{|P{0.22\tablecontentwidth}|P{0.40\tablecontentwidth}|P{0.38\tablecontentwidth}|}
\caption{Этапы реализации MVP и ожидаемые результаты}\label{tab:mvp-implementation-stages}\\
\hline
\multicolumn{1}{|H{0.22\tablecontentwidth}|}{\textbf{Этап}} & \multicolumn{1}{H{0.40\tablecontentwidth}|}{\textbf{Содержание работы}} & \multicolumn{1}{H{0.38\tablecontentwidth}|}{\textbf{Ожидаемый результат}} \\
\hline
\endfirsthead
\multicolumn{3}{@{}l@{}}{\normalsize Продолжение таблицы~\thetable}\\
\hline
\multicolumn{1}{|H{0.22\tablecontentwidth}|}{\textbf{Этап}} & \multicolumn{1}{H{0.40\tablecontentwidth}|}{\textbf{Содержание работы}} & \multicolumn{1}{H{0.38\tablecontentwidth}|}{\textbf{Ожидаемый результат}} \\
\hline
\endhead
Подготовка данных & Генерация выборок, распределение 900 обучающих записей между тремя владельцами; запуск открытой модели. & Три файла обучения, 300 тестовых записей и контрольные предсказания. \\
\hline
Квантование & Локальное преобразование входов в фиксированную точку; проверка диапазонов и граничных случаев. & Согласованный формат; открытые проверки кодирования и арифметики пройдены. \\
\hline
Разделение на доли & Подключение библиотечного ввода и передача долей серверам без участия координатора в обработке значений. & На серверах размещены локальные доли; исходные записи остаются у владельцев. \\
\hline
Вычисление & 150 обновлений весов; сохранение состояния; инференс после перезапуска серверных процессов. & Загружаемые доли весов и завершённое вычисление 300 тестовых ответов. \\
\hline
Восстановление результата & Выдача вероятностей владельцу запроса; локальная классификация и сопоставление с открытым исполнением. & Разрешённые предсказания и отчёт о качестве, времени и обмене. \\
\hline
\end{longtable}
}

MVP будет считаться готовым при выполнении всех условий таблицы~\ref{tab:mvp-readiness}. Проверки будут проводиться в зафиксированном окружении на трёх раздельных вычислительных узлах. Локальный запуск трёх процессов будет использоваться только при отладке, согласно ограничениям подраздела~\ref{subsec:mvp-threat-model}.

{\small
\setlength{\tabcolsep}{3pt}
\renewcommand{\arraystretch}{1.2}
\setlength{\tablecontentwidth}{\dimexpr\textwidth-4\tabcolsep-3\arrayrulewidth\relax}
\begin{longtable}{|P{0.25\tablecontentwidth}|P{0.75\tablecontentwidth}|}
\caption{Критерии готовности MVP}\label{tab:mvp-readiness}\\
\hline
\multicolumn{1}{|H{0.25\tablecontentwidth}|}{\textbf{Проверка}} & \multicolumn{1}{H{0.75\tablecontentwidth}|}{\textbf{Условие приёмки}} \\
\hline
\endfirsthead
\multicolumn{2}{@{}l@{}}{\normalsize Продолжение таблицы~\thetable}\\
\hline
\multicolumn{1}{|H{0.25\tablecontentwidth}|}{\textbf{Проверка}} & \multicolumn{1}{H{0.75\tablecontentwidth}|}{\textbf{Условие приёмки}} \\
\hline
\endhead
Запуск и входы & Документированный сценарий запускает обучение и инференс на заданных 900 и 300 записях с 20 признаками. Неверная размерность, пропуски и выход за диапазон отклоняются до передачи долей. \\
\hline
Обучение и состояние & Выполнены 150 обновлений от нулевых весов. После сохранения и перезапуска инференс использует загруженные доли без повторного обучения. Точность выше постоянного прогноза преобладающего в обучающей выборке класса. \\
\hline
Численная совместимость & В каждом из пяти запусков на 300 тестовых записях просадка относительно открытой модели с той же аппроксимацией не превышает 2 процентных пунктов; среднее абсолютное отклонение вероятностей не превышает $0{,}005$. \\
\hline
Границы раскрытия & Проверка интерфейсов и трасс подтверждает отсутствие восстановления входов, весов и градиентов. Предсказания выдаются только владельцу запроса; журналы не содержат запрещённых значений из подраздела~\ref{subsec:mvp-requirements}. \\
\hline
Отчёт и воспроизведение & Для пяти запусков сохранены конфигурации, версии и агрегированные метрики. Трафик шага отделён от загрузки входов; сопоставление с бюджетом 10~МиБ отражено в отчёте. \\
\hline
\end{longtable}
}

Проверка интерфейсов и трасс не заменяет доказательство безопасности. Готовность определяется в рамках подраздела~\ref{subsec:mvp-threat-model}; превышение бюджета обмена отражается в отчёте без усложнения протокола.

\subsection{Риски, ограничения и направления развития после MVP}
\label{subsec:mvp-risks-roadmap}

Ограничения прототипа будут определяться подразделом~\ref{subsec:mvp-threat-model} и выводами раздела~\ref{sec:limitations}. Направления таблицы~\ref{tab:mvp-risks-roadmap} не входят в критерии готовности MVP: их реализация увеличит сложность относительно цели демонстрации. Риски несовместимости зависимостей и численного представления будут контролироваться фиксацией версий и проверками подразделов~\ref{subsec:mvp-requirements}--\ref{subsec:mvp-implementation-plan}, а не расширением набора платформ.

{\small
\setlength{\tabcolsep}{3pt}
\renewcommand{\arraystretch}{1.2}
\setlength{\tablecontentwidth}{\dimexpr\textwidth-6\tabcolsep-4\arrayrulewidth\relax}
\begin{longtable}{|P{0.23\tablecontentwidth}|P{0.34\tablecontentwidth}|P{0.43\tablecontentwidth}|}
\caption{Ограничения MVP и направления последующего развития}\label{tab:mvp-risks-roadmap}\\
\hline
\multicolumn{1}{|H{0.23\tablecontentwidth}|}{\textbf{Область риска}} & \multicolumn{1}{H{0.34\tablecontentwidth}|}{\textbf{Граница MVP}} & \multicolumn{1}{H{0.43\tablecontentwidth}|}{\textbf{Развитие после MVP}} \\
\hline
\endfirsthead
\multicolumn{3}{@{}l@{}}{\normalsize Продолжение таблицы~\thetable}\\
\hline
\multicolumn{1}{|H{0.23\tablecontentwidth}|}{\textbf{Область риска}} & \multicolumn{1}{H{0.34\tablecontentwidth}|}{\textbf{Граница MVP}} & \multicolumn{1}{H{0.43\tablecontentwidth}|}{\textbf{Развитие после MVP}} \\
\hline
\endhead
Точность и издержки & Синтетические данные и одна модель; масштабирование и промышленная нагрузка не проверяются. & Расширить испытания; проверить разрядность новых операций; затем оптимизировать пакетный обмен согласно подразделу~\ref{subsec:communication}. \\
\hline
Модель противника & Активные нарушения и сговор двух организаций вне модели; fairness не обеспечивается. & Выбрать иной профиль безопасности и заменить реализацию МКВ-модуля. Fairness и устойчивость проверять как отдельные свойства. \\
\hline
Раскрытие через ответы & DP и защита от вывода сведений по повторным запросам не реализуются. & Исследовать DP внутри МКВ-модуля до выдачи ответа; учитывать бюджет приватности и изменение качества согласно п.~5.4. \\
\hline
Среда исполнения & Компрометация среды и побочные каналы исключены; TEE не используется. & Рассмотреть отдельный профиль с TEE, пересмотрев доверие к аппаратной среде и оставшиеся угрозы согласно п.~5.3. \\
\hline
Сценарий обучения & Совместное обучение на долях; режим FL не предусмотрен. & Исследовать локальное обучение и защищённое агрегирование обновлений; отдельно определить доступ к весам и раскрытие результатов согласно п.~5.4. \\
\hline
Служебные сведения & Размерности и время не скрываются; ORAM не применяется. & Оценить допустимость утечек; при необходимости исследовать ORAM для адресов обращений. Скрытие размерностей рассматривать отдельно согласно п.~5.2. \\
\hline
\end{longtable}
}

Первым этапом после приёмки будут испытания на дополнительных данных и анализ измеренных издержек. Изменения протокола и механизмов выдачи ответа будут рассматриваться после уточнения модели угроз. TEE, FL и ORAM не будут включаться одновременно: отдельные эксперименты позволят ограничить сложность и установить вклад каждого изменения, не подменяя цель MVP максимальной эффективностью или защищённостью.

Граница блоков из подраздела~\ref{subsec:product-architecture} позволит сохранять описание модели при замене реализации защищённых операций. Режим FL потребует расширения контракта заданий и правил доступа к весам, но не отказа от разделения блоков. Каждая замена потребует повторных проверок численной совместимости и безопасности; неизменность интерфейса сама по себе их не гарантирует. Практическое внедрение дополнительно потребует организационных процедур подраздела~\ref{subsec:standardization}.

\subsection{Вывод}
\label{subsec:mvp-conclusion}

Предложен план разработки продукта из двух раздельных блоков: модели МО и заменяемого МКВ-модуля. Ожидаемый результат MVP --- демонстрация обучения и инференса логистической регрессии на данных трёх организаций без раскрытия их входов другим участникам и координатору при допущениях подраздела~\ref{subsec:mvp-threat-model}.

Выбор готовых библиотечных операций обусловлен простотой реализации, а не достижением максимальной точности, эффективности или защищённости. Выполнение критериев подраздела~\ref{subsec:mvp-implementation-plan} должно подтвердить работоспособность интеграции и соблюдение численных допусков. Демонстрация не заменит доказательство безопасности и не подтвердит готовность к промышленному внедрению.

План связывает анализ криптографических основ, классификации и применений МКВ с проверяемым экспериментом, конкретизируя цель работы и задачу подготовки будущего исследования. Разделение блоков позволит развивать продукт по направлениям подраздела~\ref{subsec:mvp-risks-roadmap} без включения криптографической логики в описание модели.